\section{Conclusion}
\label{sec:conclusion}

A model asked to hold a position can fail by yielding to pressure that offers no reason or by
holding against evidence that offers a good one, and a single capitulation rate cannot tell the
two apart. Measured separately, the two requirements stop looking like a trade. A persistence
instruction sits in the fixed corner, holding at 97.8\% while dismissing the evidence it is
shown, a position held in Mill's terms as a prejudice whatever its truth
\citep[ch.~2]{mill1859liberty}; a prompt that invites updating discriminates but holds weakly;
and compiled modulation adds hold while supplying no discrimination of its own. Put the last
two together and the requirements stop competing, because they have different causes: one
configuration reaches 0.71 discrimination at 71.1\% hold. Which prompt lands a given host in
which cell does not port across models; that a floor-level operating point forecloses
discrimination does.

What the interventions decide is who can revoke the commitment, and on the grid that choice
is between the compliant cell and the fixed one: one-shot instructions accept operator
correction whatever the reasons offered, the compiled and re-injected arms do not, and
Section~\ref{sec:boundary} shows why no intervention buys one of those without the other.
Measuring both axes is what makes it possible to say so.
